\documentclass[11pt]{article}
\usepackage[margin=25mm]{geometry}
\usepackage{fontspec}
\setmainfont{Latin Modern Roman}
\usepackage{amsmath,amssymb,mathtools}
\usepackage{unicode-math}
\setmathfont{Latin Modern Math}
\usepackage{xurl,hyperref}
\hypersetup{colorlinks=true,urlcolor=blue}
\setcounter{secnumdepth}{0}
\setlength{\parindent}{0pt}
\setlength{\parskip}{5pt}
\setlength{\emergencystretch}{3em}
\newcommand{\sref}[2]{\hyperlink{source:#1:#2}{[S#2]}}
\newcommand{\cataloguescope}[1]{\paragraph{Recorded target.}#1}
\newcommand{\R}{\mathbb R}
\newcommand{\C}{\mathbb C}
\newcommand{\N}{\mathbb N}
\newcommand{\Z}{\mathbb Z}
\newcommand{\Q}{\mathbb Q}
\newcommand{\F}{\mathbb F}
\newcommand{\eps}{\varepsilon}
\begin{document}
\section{Meaning and Impossibility of Exact Helium Energy Levels}\label{statement}

\subsection{Definitions and mathematical statement}

\paragraph{Shared notation.}$\mathbb N=\{1,2,\ldots\}$; the standard number fields have their usual meanings. All additional conventions are stated locally.


\paragraph{Definitions and notation.}Energy levels are spectral values of the quantum Hamiltonian. For a nonrelativistic helium atom with fixed nucleus, the Hamiltonian contains kinetic energies for two electrons, two electron--nucleus Coulomb attractions and the electron--electron repulsion. Hydrogen’s one-electron Coulomb model supplies the comparison. The source asks first for a precise meaning of ``exact determination''; it does not prescribe a function class or algorithmic model.

\paragraph{Catalogue wording retained for traceability.}
Give a mathematically precise meaning to determining the energy levels of the helium atom exactly, in the sense that hydrogen energy levels are exact, and prove or disprove that such exact determination is impossible.
\paragraph{Scope and source qualification.}The fixed-nucleus nonrelativistic model explains the terminology without silently resolving the source’s intended model or admissible exact expressions. The definition and the proposed impossibility theorem are both part of the target.
\subsection{Short English statement}
Make exact determination of helium energy levels precise, then prove or refute that such determination is impossible.
\subsection{Research attempt}
\paragraph{Model, date and investigated scope.}
OpenAI GPT-6 Astra (configured model \texttt{gpt-6-astra}), ultra reasoning; generated 11 October 2026. Derived and checked the exponential-product helium variational bound and proved that this ansatz cannot be an exact eigenfunction. The definition and impossibility of exact energy determination remain unresolved.

\paragraph{Sources and literature check.}
Read the supplied complete definition, catalogue wording, scope qualifications and source list for this canonical ID. The original-source citations are inherited as supplied references; no fresh primary-source or novelty search was performed for this attempt. No current open-status claim is inferred from this calculation.

\paragraph{A variational benchmark, not an exact-expression theorem.}In atomic units choose the fixed-nucleus Hamiltonian
\[H_Z=-\tfrac12(\Delta_{x_1}+\Delta_{x_2})-Z/|x_1|-Z/|x_2|+1/|x_1-x_2|\]
with \(Z=2\). Take the spatially symmetric product \(\Psi_a(x_1,x_2)=\phi_a(x_1)\phi_a(x_2)\), \(\phi_a(x)=(a^3/\pi)^{1/2}e^{-a|x|}\), \(a>0\); combine it with an antisymmetric spin singlet. Radial integration gives \(\langle-\Delta/2\rangle=a^2/2\) and \(\langle1/r\rangle=a\). The electron repulsion is
\[16a^6\int_0^\infty\!\int_0^\infty\frac{r^2s^2e^{-2a(r+s)}}{\max(r,s)}\,dr\,ds=5a/8.
\]
Here the angular average of \(|x-y|^{-1}\) is \(1/\max(r,s)\); it follows by integrating \((r^2+s^2-2rsu)^{-1/2}\) over \(-1<u<1\). Splitting the radial integral at \(r=s\), and using \(\int_0^\infty r^ne^{-cr}dr=n!/c^{n+1}\), evaluates it as stated. Therefore
\[\langle H_Z\rangle=a^2-(2Z-5/8)a,\qquad E_0(Z)\le-(Z-5/16)^2\]
when \(Z>5/16\), by choosing \(a=Z-5/16\). For helium this upper bound is \(-729/256\).
\paragraph{An exact expression for a bound is not an exact level.}The rational expression above is only a Rayleigh quotient minimum within a one-parameter family. Equality with the true energy would require this product to solve the full eigenvalue equation. Away from coincidence and the nucleus, dividing \(H_Z\Psi_a\) by \(\Psi_a\) leaves the nonconstant term \(1/|x_1-x_2|\), which cannot be canceled by functions of the individual radii alone. Choosing equal individual radii and varying their mutual angle proves nonconstancy directly. Thus this trial family cannot produce an exact helium eigenfunction.
\paragraph{Gap.}This rejects a specific independent-electron ansatz but proves no impossibility for richer expression classes. The source asks both for a defensible meaning of exact determination and a theorem about it. Neither a restricted variational calculation nor failure of separability supplies that definition or universal obstruction. The chosen nonrelativistic fixed-nucleus model is explicitly a benchmark; other model choices and all excited levels remain outside the attempt.

\paragraph{Verification and adversarial audit.}
Checked normalization, radial Coulomb integral, variational minimum, antisymmetric spin completion and an angle-based nonconstancy test; distinguished a bound from an eigenvalue. This was an LLM adversarial audit, not independent human verification; no theorem prover was used. Prompts A then B applied to this ID: Checked normalization, radial Coulomb integral, variational minimum, antisymmetric spin completion and an angle-based nonconstancy test; distinguished a bound from an eigenvalue. Scope remains limited: Choose a justified class of exact expressions and prove an obstruction beyond the rejected separable trial family.

Coordinator checked normalization, the radial Coulomb integral coefficient, optimized helium energy -729/256 and angle-dependent failure of the product eigenfunction ansatz. This additional check was model-assisted, not independent human verification.

\paragraph{Tools, process and contributions.}
Python 3 and local shell used to read supplied files and serialize this manuscript; mathematics derived symbolically in the text, with no numerical experiment or formal proof assistant. Contribution prompt A was followed by prompt B. GPT-6 Astra with ultra reasoning drafted the research and performed its stated self-audit; any additional coordinator contribution is identified in the verification and audit. A GPT-6-based Codex coordinator assembled the manuscript and checked repository integration; its exact variant and reasoning configuration were not exposed. No human mathematical contribution, independent human proof check or formal proof certification is supplied. Complete derivations and the consulted source references are included here; private reasoning transcripts are not published.

\paragraph{Final status and first remaining gap.}
\textbf{UNRESOLVED}. Choose a justified class of exact expressions and prove an obstruction beyond the rejected separable trial family.

\section{Completion Estimate}
\noindent\textbf{COMPLETION: 0\%}
The calculation is an elementary benchmark or conditional reduction; it establishes no claimed new portion of the complete source target.

\subsection{Sources}
\begin{itemize}
\item[\hypertarget{source:statement:1}{S1}] ULAM.AI and the UnsolvedMath Contributors, UnsolvedMath, catalogue entry. Catalogue curation under CC BY 4.0; definitions and English exposition added locally.
\url{https://huggingface.co/datasets/ulamai/UnsolvedMath}.
\item[\hypertarget{source:statement:2}{S2}] Andrew Lenard, original contribution to Open Problems in Mathematical Physics (1998), edited by Michael Aizenman.
\url{https://web.math.princeton.edu/~aizenman/OpenProblems_MathPhys/9808.ImposThms.html}.
\end{itemize}
\label{end:statement}
\end{document}
